\documentclass[11pt,a4paper]{article}

\usepackage[margin=1in]{geometry}
\usepackage{amsmath,amssymb,amsthm}
\usepackage{mathtools}
\usepackage{algorithm}
\usepackage{algpseudocode}
\usepackage{booktabs}
\usepackage{microtype}
\usepackage[hidelinks]{hyperref}
\usepackage{natbib}
\usepackage{url}

\newtheorem{conjecture}{Conjecture}
\newtheorem{problem}{Open Problem}
\newtheorem{condition}{Sufficient Condition}

\DeclareMathOperator*{\argmax}{arg\,max}
\DeclareMathOperator*{\argmin}{arg\,min}

\title{A Joint Minimization Framework for Transformer Interpretability\\
       \large\textit{Information Crystallization}}
\author{Saqib Nazir Bhat\thanks{Independent research. Contact: \texttt{saqibnazirbhat3@gmail.com}. Repository: \url{https://github.com/Saqibnazirbhat/information-crystallization}}}
\date{March 2026}

\begin{document}
\maketitle

\begin{abstract}
\noindent
Transformer interpretability is usually posed as: identify the parameters and
inputs responsible for a model's behaviour. Existing methods work top-down,
searching the combinatorial space by ablation from the full model. I propose a
different formulation. Given an observed output token $\hat{x}$, find the
minimal $S \subseteq [d]$ and $T \subseteq [n]$ such that the argmax over the
restricted model still equals $\hat{x}$. A single observation simplifies the
problem: softmax does not affect argmax. Preserving the predicted token requires
only that the logit margin $\delta$ remains positive, which is strictly weaker
than matching the full distribution. Under this relaxation I introduce
\emph{Information Crystallization}, a bottom-up $O(k\cdot d)$ algorithm
(seed, greedy growth, prune) that finds local minima of $(S, T)$ jointly. I
decompose $S$ into a final-block component $S_L$ and a residual-stream support
$S_{\text{supp}}$, connect the sparsity conjecture to the Lottery Ticket
Hypothesis, and identify the $T$-dependence of $S^*$ as the central open
problem. Preliminary results from the seed phase on GPT-2 small (N$=$17,
stratified by query type) refute the suffix-only heuristic for input-position
seeding; only $11.8\%$ of seed positions land on the last token, and
long-range sequences cluster near the subject noun (mean distance 7.67 tokens
from the end). The parameter-seed result is inconclusive on GPT-2 because its
input embedding and unembedding projection share one parameter tensor;
decisive evaluation of the $S_L \cup S_{\text{supp}}$ decomposition requires
untied architectures. Growth and prune phases are unimplemented.
\end{abstract}

\section{Introduction}

Mechanistic interpretability of large language models has become one of the
central open problems in AI safety. Prior work treats parameter attribution and
input attribution as separate questions: circuit discovery~\citep{elhage2021transformer},
activation patching, causal tracing~\citep{meng2022locating}, and input
rationales~\citep{lei2016rationalizing,bau2020identifying} all prune components
starting from the full model. The search space is $O(2^d)$ for parameters
($d \sim 10^{11}$) and $O(2^n)$ for inputs. Practical methods rely on
heuristics to make this tractable.

This paper takes a different route. I observe that the practical interpretability
question---``why did the model produce \emph{this} token?''---does not require
reproducing the model's full output distribution. It requires only that the
restricted model still predicts the same token. That is a weaker condition, and
it is the key lever for the rest of the paper.

\paragraph{Contributions.}
\begin{enumerate}
  \item A joint formulation of parameter and input minimization with a single
        preservation condition (Section~\ref{sec:problem}).
  \item The \emph{softmax relaxation}: replacing distribution preservation with
        a margin inequality (Section~\ref{sec:relaxation}).
  \item \emph{Information Crystallization}, a bottom-up algorithm with seed,
        growth, and prune phases (Section~\ref{sec:algorithm}).
  \item A residual-stream decomposition $S = S_L \cup S_{\text{supp}}$
        (Section~\ref{sec:s-decomp}) and a formal statement of the coupling
        problem (Section~\ref{sec:coupling}).
  \item An empirical protocol on GPT-2 together with preliminary seed-phase
        results on $N{=}17$ stratified sequences (Section~\ref{sec:protocol}).
\end{enumerate}

\paragraph{Relation to prior work.} Information Crystallization is most
closely related to attribution patching~\citep{syed2023attribution}, which
also uses gradient approximations to circumvent the combinatorial cost of
full ablation. The critical difference is direction: attribution patching
scores every parameter once by its linearized ablation effect; Information
Crystallization iteratively grows a set from a single seed. The two methods
are complementary---attribution patching produces a dense importance map;
Information Crystallization produces a minimal sufficient subset. Automated
Circuit Discovery (ACDC)~\citep{conmy2023automated} is the canonical
top-down baseline we contrast with. The bottom-up construction also echoes
inference-time pruning methods such as SparseGPT~\citep{frantar2023sparsegpt}
and Wanda~\citep{sun2023simple}, though those operate globally on the
weight distribution rather than per-prediction. The gradient-attribution
step builds on Integrated Gradients~\citep{sundararajan2017axiomatic}.

\section{Problem Formulation}
\label{sec:problem}

Let $V$ be a finite vocabulary and $X = (x_1, \ldots, x_n) \in V^n$ an input
sequence. A language model parameterized by $\theta \in \mathbb{R}^d$ defines
\begin{equation}
  P_\theta(x_{n+1} \mid x_1, \ldots, x_n) = \mathrm{softmax}\!\big(f_\theta(X)\big),
\end{equation}
where $f_\theta : V^n \to \mathbb{R}^{|V|}$ is a composition of $L$ non-linear
transformations:
\begin{equation}
  f_\theta = f_\theta^{(L)} \circ f_\theta^{(L-1)} \circ \cdots \circ f_\theta^{(1)}.
\end{equation}
The predicted token is
\begin{equation}
  \hat{x} = \argmax_{v \in V} P_\theta(v \mid X).
\end{equation}

\paragraph{Problem.} Given $\hat{x}$, find minimal $S \subseteq [d]$ and
$T \subseteq [n]$ such that
\begin{equation}
  \argmax_{v \in V} P_{\theta_S}(v \mid X_T) = \hat{x},
\end{equation}
where $\theta_S$ denotes $\theta$ restricted to indices $S$ and $X_T$ denotes
$X$ restricted to positions $T$.

\paragraph{Operationalization of $\theta_S$.} Throughout this paper I take
zero-ablation: entries outside $S$ are set to zero. This is the simplest
commitment and is what Algorithm~\ref{alg:crystallize} evaluates. Alternative
operationalizations---mean-ablation over a reference distribution, or
activation-patching from a paired input---are compatible with the framework
and in general yield different minimal $(S, T)$. The choice matters empirically:
zero-ablation is conservative (it over-attributes because the network sees
degenerate inputs it was never trained on), while mean-ablation is more
permissive. Future work will compare all three.

\paragraph{Tokenization.} $T$ is defined at the token level of the model's
tokenizer (BPE for GPT-2). A BPE-split word therefore occupies multiple
positions and may be partially selected. Claims about word-level salience
must be recovered post-hoc from token-level $T$.

In general, this problem is computationally intractable: $f_\theta$ is
non-linear and non-invertible, and $d \sim 10^{11}$ for current
frontier-scale models. Section~\ref{sec:relaxation} shows that a natural
weakening of the preservation condition makes the problem substantially more
tractable.

\section{The Softmax Relaxation}
\label{sec:relaxation}

\subsection{Margin definition}

The logit margin is the gap between the winning logit and its nearest
competitor:
\begin{equation}
  \delta \;=\; f_\theta(X)[\hat{x}] \;-\; \max_{v \neq \hat{x}} f_\theta(X)[v].
\end{equation}
The argmax condition $\argmax_v f_\theta(X)[v] = \hat{x}$ is equivalent to
$\delta > 0$. Since softmax is a strictly monotone transformation, preserving
the argmax requires only preserving the \emph{sign} of $\delta$. No particular
probability value needs to match.

This is the key relaxation. A distribution-preservation condition
becomes a margin-preservation condition. The former is strict; the
latter admits a buffer zone of width $\delta$ around the decision boundary
within which perturbations do not change the argmax.

\subsection{The $\delta$-buffer argument}

Most parameters in a language model do not determine the identity of the
winning token. They refine probabilities of losing tokens, smooth the output
distribution, or encode features relevant to positions outside $T$. A parameter
$\theta_i$ is irrelevant to $\hat{x}$ if perturbing it changes $\delta$ by less
than $\delta$ itself. The set of such parameters is large precisely because
$\delta > 0$ creates a buffer around the decision boundary. This motivates the
central \emph{Sparse Explanation Hypothesis}: the true minimal $(S, T)$ is
small relative to $(d, n)$.

\subsection{Sufficient condition}

The following condition, if proved, converts the conjectures below into
theorems.

\begin{condition}\label{cond:lipschitz}
  For all $\theta_i \notin S$,
  \begin{equation}
    \bigl| \mathrm{logit}_{\hat{x}}(\theta_S, X_T) - \mathrm{logit}_{\hat{x}}(\theta, X) \bigr| < \delta.
  \end{equation}
\end{condition}

A proof via Lipschitz continuity of $f_\theta$ or a first-order perturbation
bound would suffice. The difficulty is the non-linearity of $f_\theta$, which
resists tight analytical bounds. A first-order Taylor expansion gives
\[
  \Delta \mathrm{logit}_{\hat{x}} \;\approx\; \nabla_\theta \mathrm{logit}_{\hat{x}}^\top \cdot \Delta \theta,
\]
suggesting that parameters with small gradient magnitude are natural candidates
for the complement $S^c$.

\section{Information Crystallization}
\label{sec:algorithm}

Top-down methods search the space $O(2^d)$ by ablating from $|S| = d$. I invert
the direction of search: start with $|S| = 1$ and grow. If the true minimal
$(S, T)$ is small, growing finds it in time proportional to its size rather
than the model's size.

\subsection{Algorithm}

\begin{algorithm}[H]
\caption{Information Crystallization}
\label{alg:crystallize}
\begin{algorithmic}[1]
\State $g \gets \nabla_\theta \log P_\theta(\hat{x} \mid X)$
\State $S \gets \{\argmax_j |g_j|\}$, \quad $T \gets \{\argmax_t \|\partial f_\theta / \partial x_t\|\}$
\Comment{seed phase}
\Repeat
    \State \textit{Forward pass on restricted subnetwork:} evaluate $\argmax_v P_{\theta_S}(v \mid X_T)$.
    \If{it equals $\hat{x}$} \textbf{break} \EndIf
    \State $\hat{g} \gets \nabla_\theta \log P_{\theta_S}(\hat{x} \mid X_T)$
    \Comment{gradient w.r.t.\ \emph{all} $\theta$, not only $\theta_S$}
    \State $S \gets S \cup \{\argmax_{j \notin S} |\hat{g}_j|\}$
\Until{$\argmax_v P_{\theta_S}(v \mid X_T) = \hat{x}$}
\Comment{growth phase}
\ForAll{$s \in S$}
    \Comment{crystallization phase}
    \If{$\argmax_v P_{\theta_{S \setminus \{s\}}}(v \mid X_T) = \hat{x}$}
        $S \gets S \setminus \{s\}$
    \EndIf
\EndFor
\If{argmax no longer preserved after pruning} \textbf{goto line 3} \Comment{verification loop}
\EndIf
\State \Return $(S, T)$
\end{algorithmic}
\end{algorithm}

\paragraph{Note on line 6.} The gradient $\hat{g}$ is taken with respect to
the full parameter vector $\theta$, not $\theta_S$. Only the forward pass
uses the restricted subnetwork. This is essential: to rank candidates
$j \notin S$, the gradient component $\hat{g}_j$ must be defined for all
$j \in [d]$.

\subsection{Complexity}

The seed phase costs one full backward pass, $O(d)$. Each growth step costs a
forward--backward pass through a restricted subnetwork of current size $|S|$
plus a global ranking of candidates at cost $O(d)$. Summing over $k$ growth
steps gives $O(k \cdot d + \sum_{i=1}^{k} i) = O(k \cdot d + k^2)$, which
simplifies to $O(k \cdot d)$ when $k \ll \sqrt{d}$ (the regime of practical
interest under the Sparse Explanation Hypothesis). The prune phase costs
$O(k)$ forward passes. Compared to exhaustive search $O(2^d)$ or top-down
pruning $O(d^2)$, this is tractable.

\subsection{Why bottom-up}

Three properties distinguish this approach:
\begin{itemize}
  \item \textbf{Direction of search.} Existing methods (activation patching,
        structured pruning, circuit discovery) start from the full model and
        remove components. This algorithm starts from empty and adds.
  \item \textbf{Gradient-guided growth.} Ranking $d$ candidates at each step
        would cost $O(d^2)$ in total. Using the gradient of the restricted
        subnetwork as a ranking heuristic reduces this to $O(d)$ per step.
  \item \textbf{Joint optimization.} $S$ and $T$ grow simultaneously,
        capturing interactions between which parameters matter and which
        inputs matter.
\end{itemize}

\section{Minimizing $T$: Input Position Selection}
\label{sec:t-min}

Let $\alpha^{h,\ell}(t)$ denote the attention weight on position $t$ from head
$h$ at layer $\ell$. Define the cumulative salience
\begin{equation}
  \sigma(t) \;=\; \sum_{h,\ell} \alpha^{h,\ell}(t).
\end{equation}
For a threshold $\varepsilon$ chosen so that positions with $\sigma(t) <
\varepsilon$ contribute less than $\delta/2$ to the winning logit, define
\begin{equation}
  T_{\text{sal}} \;=\; \{\, t \in [n] : \sigma(t) > \varepsilon \,\}.
\end{equation}

\begin{conjecture}\label{conj:suffix}
  $T = \{x_{n-k}, \ldots, x_n\} \cup T_{\text{sal}}$ with $k \ll n$, where the
  suffix captures recency and $T_{\text{sal}}$ captures earlier positions with
  high semantic weight.
\end{conjecture}

\paragraph{Failure modes.} Transformers use global attention, so a position-1
token with high $\sigma$ can dominate the logit landscape. The suffix
approximation breaks in three concrete cases:
\begin{itemize}
  \item Sentence-initial negation: ``Despite everything, the answer is still\ldots''
        Position 1 reshapes the entire logit distribution, overriding suffix
        recency.
  \item Long-range syntactic dependencies: the subject introduced at positions
        1--3 may determine verb form 20 or more tokens later.
  \item Named entities: in ``Harry Potter picked up his\ldots'', the entity at
        positions 1--2 causes the model to prefer ``wand'' over ``phone''. The
        suffix alone is insufficient.
\end{itemize}
The corrected claim is: $T$ is identified by attention weight concentration,
not by positional recency. The suffix approximation is a useful empirical
heuristic but not a structural guarantee. I formalize the corrected claim:

\begin{conjecture}[attention-concentration identification]\label{conj:attention}
  For a minimal $T$ that preserves $\hat{x}$ under the restricted model,
  every $t \in T$ satisfies $\sigma(t) > \varepsilon$ for a threshold
  $\varepsilon$ determined by the decision boundary. In particular,
  Conjecture~\ref{conj:suffix} is superseded: the suffix $\{x_{n-k}, \ldots, x_n\}$
  is not structurally guaranteed to lie in $T$.
\end{conjecture}

Preliminary seed-phase data on GPT-2 small (Section~\ref{sec:protocol})
is consistent with Conjecture~\ref{conj:attention}: the seed input position
$t^*$ lies on the final token in only $11.8\%$ of sequences, and tracks the
subject noun on long-range dependency queries. Conjecture~\ref{conj:suffix}
remains in the paper as a deprecated reference point.

\section{Minimizing $S$: Parameter Selection}
\label{sec:s-decomp}

\subsection{Residual stream decomposition}

I conjecture that $S$ decomposes as
\begin{equation}
  S \;=\; S_L \cup S_{\text{supp}},
\end{equation}
where $S_L$ consists of parameters in the final transformer block that
directly govern $\delta$ (specifically: the unembedding projection weights for
$\hat{x}$ and its nearest competitor, the attention heads in the last one or
two layers that activate most strongly on $T$, and a small set of MLP neurons
in the final block), and $S_{\text{supp}}$ consists of earlier-layer parameters
whose representations are consumed by $S_L$ via the residual
stream~\citep{elhage2021transformer}. The contribution of $S_{\text{supp}}$
must satisfy
\begin{equation}
  \bigl| \delta(S_L) - \delta(S_L \cup S_{\text{supp}}) \bigr| < \delta / 2.
\end{equation}

\paragraph{Weight-tying caveat.} Equation~(8) presumes that the unembedding
projection and the input token-embedding matrix are distinct parameter
tensors. This presumption fails on GPT-2 and many other pretrained models,
which tie \texttt{lm\_head.weight} to \texttt{transformer.wte.weight}.
On weight-tied models the $S_L \cup S_{\text{supp}}$ decomposition is
ill-posed as stated: gradient flow through the unembedding projection and
through the input embedding is indistinguishable at the level of the
parameter tensor. The seed-phase experiment in Section~\ref{sec:protocol}
exposes exactly this. A clean empirical test of the decomposition therefore
requires either (a) an architecture with untied embeddings
(e.g., Pythia, OPT), or (b) an experimental design that separately attributes
gradient mass to the embedding-path and unembedding-path contributions to a
shared tensor. Both are future work.

\subsection{Lottery Ticket connection}

\citet{frankle2019lottery} show that dense networks contain sparse subnetworks
---``winning tickets''---that match the full network's performance when
trained in isolation. The conjecture here is the inference-time analogue: for
a fixed input $X$ and prediction $\hat{x}$, there exists a sparse subnetwork
that preserves the argmax without retraining.

\subsection{Middle-layer localization}

\citet{meng2022locating} show via causal tracing that factual associations in
GPT-class models localize in middle-layer MLP modules. This challenges the
assumption that $S_L$ alone suffices for factual recall queries. For ``The
capital of France is\ldots'' the decisive parameters may live in middle
layers, meaning $S_{\text{supp}}$ could be substantially larger for factual
than for syntactic continuation tasks. Any empirical protocol must stratify by
query type.

\section{Joint Minimization: The Coupling Problem}
\label{sec:coupling}

$S$ and $T$ are not independently minimizable. When $T$ changes, the embeddings
fed to layer 1 change. That perturbation propagates through all $L$ layers via
the residual stream, producing a different hidden state at layer $L$. The set
of parameters in $S_L$ that are load-bearing for $\delta$ therefore depends on
$T$:
\begin{equation}
  S^*(T) \;=\; \argmin_S |S| \quad \text{s.t.} \quad \argmax_v P_{\theta_S}(v \mid X_T) = \hat{x}.
\end{equation}
Treating $S$ and $T$ as independent is a useful approximation but its error is
only small when $S_{\text{supp}}$ is small relative to $\delta$.

\begin{problem}
  Characterize the function $T \mapsto S^*(T)$. Two concrete sub-questions:
  (i) bound $|S^*(T) - S^*(\theta, X)|$ as a function of $|T|$; (ii) determine
  whether the greedy decomposition in Algorithm~\ref{alg:crystallize} is
  within a constant factor of the true joint minimum.
\end{problem}

\paragraph{Greedy decomposition as upper bound.} Despite the coupling,
Algorithm~\ref{alg:crystallize} yields a computable upper bound on $|S| + |T|$:
identify $T$ from cumulative attention weights (Section~\ref{sec:t-min}), then
minimize $S$ conditional on $T$ via gradient-sensitivity scoring, then check
whether the coupling error is small. This is suboptimal in general but
empirically verifiable.

\section{Empirical Protocol}
\label{sec:protocol}

GPT-2~\citep{radford2019language} is suitable for initial verification because
its full parameter set and attention weights are publicly accessible. I
propose the following protocol and report preliminary seed-phase results
in Section~\ref{sec:preliminary}.

\subsection{Protocol}

\begin{enumerate}
  \item \textbf{Sampling.} Draw 1{,}000 input sequences stratified by query
        type: syntactic continuation, factual recall, negation-heavy, and
        long-range syntactic dependency.
  \item \textbf{Baseline.} Record $\hat{x}$ for each sequence using the full
        model $\theta$.
  \item \textbf{$T$ identification.} Greedily remove input positions in
        ascending order of $\sigma(t)$. Record the minimal $T$ at which
        $\hat{x}$ first flips.
  \item \textbf{$S$ identification.} Greedily ablate parameters in ascending
        order of gradient-sensitivity score. Record the minimal $S$ at which
        $\hat{x}$ first flips.
  \item \textbf{Size measurement.} Report $|S|/d$ and $|T|/n$ distributions
        per query type (means, medians, 90th percentiles).
  \item \textbf{Coupling test.} For each $T$ from Step 3, recompute $S^*(T)$
        and compare with $S^*(\theta, X)$. Small discrepancy empirically
        validates the greedy decomposition.
\end{enumerate}

\subsection{Preliminary seed-phase results}
\label{sec:preliminary}

I have run the seed phase (Step 1 of Algorithm~\ref{alg:crystallize}) on
$N = 17$ hand-curated sequences stratified into four query types
(syntactic continuation, factual recall, negation-heavy, long-range syntactic
dependency). Sample size is a sanity check, not a statistical study.

\paragraph{Input seed $t^*$.} Across the 17 sequences, only $11.8\%$
($2/17$) had $t^*$ on the final input token. Mean distance of $t^*$ from
the sequence end was $3.88$ tokens. Stratified:

\begin{center}
\begin{tabular}{lc}
\toprule
Query type & Mean distance of $t^*$ from end \\
\midrule
factual ($n{=}5$)    & 2.00 \\
syntactic ($n{=}5$)  & 2.80 \\
negation ($n{=}4$)   & 4.75 \\
long-range ($n{=}3$) & \textbf{7.67} \\
\bottomrule
\end{tabular}
\end{center}

This is evidence---subject to the obvious caveat of small $n$, especially
$n{=}3$ in the long-range bucket---consistent with
Conjecture~\ref{conj:attention} (attention concentration, not suffix recency)
and inconsistent with Conjecture~\ref{conj:suffix} (suffix only). On
long-range sequences, $t^*$ lands on the subject noun rather than on any
suffix position.

\paragraph{Parameter seed $j^*$ (inconclusive due to weight-tying).}
In $17/17$ sequences, the seed parameter tensor was
\texttt{transformer.wte.weight}. This is uninformative for testing the
$S_L \cup S_{\text{supp}}$ decomposition as discussed in
Section~\ref{sec:s-decomp}: on GPT-2 this single tensor serves as both the
input embedding and the unembedding projection. The result is consistent
with the $S_L$ conjecture (unembedding $\in S_L$) but does not adjudicate
between $S_L$ and the input embedding as the dominant direct-gradient path.

\paragraph{A stronger reading.} The concentration of the direct gradient
of $\log P(\hat{x} \mid X)$ on a single tensor suggests that the direct
(unembedding-path) contribution to the winning logit dominates the
residual-stream-mediated contribution in the seed step. If this holds
more generally, $S_L$ is the empirically dominant component and
$S_{\text{supp}}$ functions as a refinement, consistent with the paper's
framing. An untied-architecture experiment is required to confirm.

\paragraph{Implementation.} Available in the repository at
\texttt{experiments/seed\_phase.ipynb} with full results in
\texttt{experiments/results.md}.

\section{Limitations}

\begin{itemize}
  \item \textbf{The suffix approximation is not a theorem.} $T$ is identified
        by attention analysis per input. No structural guarantee holds.
  \item \textbf{$|S_{\text{supp}}|$ is unbounded.} The inequality
        $|\delta(S_L) - \delta(S_L \cup S_{\text{supp}})| < \delta/2$ is a
        requirement, not a proved result. For factual queries it may be
        violated.
  \item \textbf{Greedy decomposition is suboptimal.} The true joint minimum
        of $(S, T)$ may be strictly smaller than the greedy bound. The gap is
        uncharacterized.
  \item \textbf{Condition~\ref{cond:lipschitz} is unproved.} All conjectures
        depend on a Lipschitz or first-order perturbation bound on $f_\theta$
        that has not been derived.
  \item \textbf{Scale.} Verification on GPT-2 does not guarantee
        generalization to models with $d \sim 10^{11}$. Distributed
        representations in large models may yield dense $(S, T)$.
  \item \textbf{Only the seed phase has been executed.} Growth, prune, and
        verification phases of Algorithm~\ref{alg:crystallize} remain
        unimplemented. Claims about the full minimal $(S, T)$ are therefore
        untested.
  \item \textbf{Weight-tying confound on GPT-2.} The preliminary $j^*$ result
        cannot cleanly test the $S_L \cup S_{\text{supp}}$ decomposition
        because GPT-2 ties input embedding and unembedding into a single
        parameter tensor. An untied-architecture replication (e.g., Pythia,
        OPT) is required.
  \item \textbf{Sample size is not statistically significant.} The
        $N{=}17$ seed-phase sample is a feasibility check. Scaling to
        $N{\geq}1{,}000$ per query type with a permutation test on
        positional distribution is the next quantitative milestone.
\end{itemize}

\section{Conclusion}

The standard interpretability problem asks for the parameters and inputs that
explain a model's behaviour. Posed as distribution preservation, the problem
is intractable. Posed as margin preservation it admits substantially sparser
solutions, reachable via a bottom-up $O(k \cdot d)$ algorithm.

The central open problem is the coupling between $S$ and $T$ through the
residual stream. A formal perturbation bound on $f_\theta$ would turn the
conjectures here into theorems. Preliminary seed-phase results on GPT-2
small are reported in Section~\ref{sec:preliminary} and
\texttt{experiments/results.md}; growth and prune phases, and replication
on an untied architecture, are the immediate next steps.

\bibliographystyle{plainnat}
\bibliography{references}

\end{document}
